% 2026-11-16 站主問：「這案子是什麼案子我都不知道」（指第 4 部所引 Bush v. Gore、第 5 部所引 Cooper v. Aaron）。
% 兩案事實與引句自 Cornell LII 判決全文核對（2026-11-16）；LII 無 U.S. Reports 頁碼標記，故引句不附頁，唯 Cooper 18 頁之至上宣告依 Law 784 頁註 198 所引。
% Stevens 不同意見書末句未取得原文，正文不引。
\chapter{所引兩案之背景}
\label{ch:80_cases}

\begin{qaquestion}
第 4 部引 Bush v. Gore\ix{cases}{02@美國!01@聯邦最高法院!2000Bush v. Gore（2000）@Bush v. Gore（2000）}，第 5 部引 Cooper v. Aaron\ix{cases}{02@美國!01@聯邦最高法院!1958Cooper v. Aaron（1958）@Cooper v. Aaron（1958）}。
\begin{enumerate}
\tightlist
\item 兩案的事實、爭點與判決為何？
\item Fallon、Law、Alexander 與 Schauer 各自在什麼脈絡下引用？
\end{enumerate}
\end{qaquestion}

兩案均為美國聯邦最高法院的裁判。Bush v. Gore 終止佛羅里達州 2000 年總統選舉的重新計票；Cooper v. Aaron 駁回阿肯色州小岩城學區延後廢除種族隔離的聲請，並宣告聯邦司法機關於憲法解釋居於至上地位。

\section{Bush v. Gore, 531 U.S. 98 (2000)}
\label{sec:80_cases-bush}

\begin{enumerate}
\tightlist
\item
  2000 年總統選舉由共和黨候選人 George W. Bush 與民主黨候選人 Al Gore 角逐，勝負取決於佛羅里達州的選舉人票。兩人在該州的票數差距極小，Gore 提起選舉爭訟。2000 年 12 月 8 日，佛羅里達州最高法院判命所有尚未以人工計算「undervotes」〔機器未讀出選擇之選票〕的郡，一律以人工重新計票。12 月 9 日，聯邦最高法院准許 Bush 停止執行的聲請，並受理上訴。\footcite{bushvgore2000}
\item
  12 月 12 日，聯邦最高法院以不具名的法院意見（per curiam）認為，各郡判斷選民意思的標準不一，甚至同一郡內各計票小組的標準也不一致，重新計票欠缺確保其平等適用的具體標準，違反增修條文第 14 條\ix{statutes}{02@美國!01@聯邦憲法!114@增修條文第 14 條}的平等保護條款。法院意見並載「Seven Justices of the Court agree that there are constitutional problems with the recount」〔本院七位大法官均同意此一重新計票有憲法上的問題〕，歧見只在救濟。\footcite{bushvgore2000}
\item
  法院意見認為，依 3 U.S.C. § 5\ix{statutes}{02@美國!02@聯邦法典第 3 編!005@第 5 條}，以終局選定選舉人為目的之爭訟應於 12 月 12 日前完成，佛州法院已無時間訂定合於平等保護的標準，遂廢棄佛州最高法院命續行重新計票的判決。這一部分為五比四。法院意見另載「Our consideration is limited to the present circumstances」〔本院之考量限於本案情形〕。Rehnquist\ix{persons}{rehnquist@Rehnquist, William H.} 首席大法官另提協同意見書，Scalia 與 Thomas 加入；Stevens、Souter、Ginsburg、Breyer 四位大法官各提不同意見書。判決次日 Gore 承認敗選。\footcite{bushvgore2000}
\item
  Fallon 20--21 頁以本案說明正當性並非自明的事實：許多贊同 Roe 的人主張本案不具正當性，也有許多論者預料 Roe 與本案會使最高法院失去更一般意義的正當性。\footcite[20--21]{fallon2018}他其後指出，批評者的指控多屬濫用裁量（依黨派動機判決）與逾越權限（依無權考慮的理由判決），係法律意義的正當性。Law 782--783 頁則說，判決之前人們可以合理懷疑最高法院能否違反多數選民的意願決定總統選舉的結果；判決作成而各方服從，使人普遍相信類似案件也會被服從，法院的權力因此增加。\footcite[782--783]{law2009}兩人的比較見第 4 部。
\end{enumerate}

\section{Cooper v. Aaron, 358 U.S. 1 (1958)}
\label{sec:80_cases-cooper}

\begin{enumerate}
\tightlist
\item
  1954 年 Brown v. Board of Education\ix{cases}{02@美國!01@聯邦最高法院!1954Brown v. Board of Education（1954）@Brown v. Board of Education（1954）} 判公立學校的種族隔離違憲；1955 年 Brown II 要求各地「with all deliberate speed」〔以審慎的速度〕廢除隔離，執行交由下級法院監督。Law 783 頁說，Brown II 的寬鬆措辭既不能防止規避，也不能促成服從。\footcite[783]{law2009}
\item
  小岩城學區依其廢除隔離計畫，安排黑人學生於 1957 年秋季進入中央高中。1957 年 9 月 2 日，阿肯色州長 Faubus 派州國民兵進駐中央高中，將校園列為禁區，阻止黑人學生入學；其後總統派遣聯邦部隊進駐中央高中，學生在部隊保護下入學。1958 年 2 月，學區以社會動盪為由，聲請將廢除隔離計畫延後，後來說明為兩年半。聯邦地方法院准許，第八巡迴上訴法院廢棄。\footcite{cooper1958}
\item
  聲請人為小岩城學區董事會成員 William G. Cooper 等人與學區督學 Virgil T. Blossom；相對人為 John Aaron 等黑人學童。\footcite{cooper1958}
\item
  聯邦最高法院於 1958 年 9 月 11 日開庭辯論，次日以法院意見維持上訴法院的判決，9 月 29 日公布理由。理由載「The constitutional rights of respondents are not to be sacrificed or yielded to the violence and disorder」〔相對人的憲法權利，不因暴力與騷亂而犧牲或讓步〕；州的官員受 Brown 案拘束；並依聯邦憲法第 6 條\ix{statutes}{02@美國!01@聯邦憲法!006@第 6 條}的至上條款與 Marbury v. Madison\ix{cases}{02@美國!01@聯邦最高法院!1803Marbury v. Madison（1803）@Marbury v. Madison（1803）}，宣告「the federal judiciary is supreme in the exposition of the law of the Constitution」〔聯邦司法機關於憲法之闡釋居於至上地位〕。理由書由九位大法官逐一具名；Frankfurter 另提協同意見書，開頭說明他在全面參與共同意見之外，另以個人名義論述。\footcite{cooper1958}
\item
  Law 784--785 頁敘述，Brown 之後數年，全國都有理由懷疑聯邦政府是否會執行學校廢除隔離的判決；艾森豪派遣第 101 空降師進駐小岩城之後，法院才重新主張其權力，Cooper 的至上宣告是在服從已有保障時作成的。\footcite[784--785 \& nn.196--198]{law2009}Alexander 與 Schauer 457 頁則把本案的至上宣告當作有待證成的規範主張。\footcite[457]{alexander2000}兩說的比較見第 5 部。
\end{enumerate}
